\documentclass[11pt,a4paper]{article}
\newcommand{\intestazione}{Dimero e trimero ad anello --- punti di lavoro e ansatz}
% =====================================================================
%  Preambolo comune ai documenti sul dimero (serie "che cosa abbiamo fatto")
%  Incluso con % =====================================================================
%  Preambolo comune ai documenti sul dimero (serie "che cosa abbiamo fatto")
%  Incluso con % =====================================================================
%  Preambolo comune ai documenti sul dimero (serie "che cosa abbiamo fatto")
%  Incluso con \input{preambolo_dimero} subito dopo \documentclass.
% =====================================================================
\usepackage[T1]{fontenc}
\usepackage[utf8]{inputenc}
\usepackage[main=italian,provide=*]{babel}
\usepackage{amsmath,amssymb,mathtools}
\usepackage{graphicx}
\usepackage{booktabs}
\usepackage{colortbl}
\usepackage{xcolor}
\usepackage{geometry}
\usepackage{placeins}
\usepackage{caption}
\usepackage{enumitem}
\usepackage{fancyhdr}
\usepackage{needspace}
\usepackage{tikz}
\usetikzlibrary{arrows.meta,positioning,shapes.geometric,calc,fit,backgrounds}
\usepackage{hyperref}
\hypersetup{colorlinks=true,linkcolor=blu,citecolor=blu,urlcolor=blu}

\geometry{margin=2.5cm}

% --- colori coerenti con quelli delle figure (palette Okabe-Ito)
\definecolor{blu}{HTML}{0072B2}
\definecolor{arancio}{HTML}{E69F00}
\definecolor{verdes}{HTML}{009E73}
\definecolor{rossos}{HTML}{D55E00}
\definecolor{violas}{HTML}{CC79A7}
\definecolor{grigetto}{HTML}{E8E8E8}

\captionsetup{font=small,labelfont=bf,width=0.92\textwidth}

\pagestyle{fancy}
\fancyhf{}
\fancyhead[L]{\small\itshape\intestazione}
\fancyhead[R]{\small\thepage}
\renewcommand{\headrulewidth}{0.3pt}

% --- notazione
\newcommand{\ket}[1]{\lvert #1\rangle}
\newcommand{\bra}[1]{\langle #1 \rvert}
\newcommand{\media}[1]{\langle #1 \rangle}
\newcommand{\vs}[1]{\mathbf{s}_{#1}}

% --- riquadro "in una riga": il messaggio da portare a casa
\newcommand{\messaggio}[1]{%
  \begin{center}
  \begin{minipage}{0.93\textwidth}
  \setlength{\fboxsep}{9pt}%
  \colorbox{grigetto}{\begin{minipage}{\dimexpr\textwidth-2\fboxsep}
  \small\textbf{In una riga.}\enspace #1
  \end{minipage}}
  \end{minipage}
  \end{center}
}

% --- riquadro "come lo sappiamo": la verifica numerica dietro un'affermazione
\newcommand{\verifica}[1]{%
  \begin{center}
  \begin{minipage}{0.93\textwidth}
  \small\textcolor{verdes}{\rule{2.2em}{1.1pt}}\enspace
  \textcolor{verdes}{\textbf{Come lo sappiamo.}}\enspace #1
  \end{minipage}
  \end{center}
  \medskip
}
 subito dopo \documentclass.
% =====================================================================
\usepackage[T1]{fontenc}
\usepackage[utf8]{inputenc}
\usepackage[main=italian,provide=*]{babel}
\usepackage{amsmath,amssymb,mathtools}
\usepackage{graphicx}
\usepackage{booktabs}
\usepackage{colortbl}
\usepackage{xcolor}
\usepackage{geometry}
\usepackage{placeins}
\usepackage{caption}
\usepackage{enumitem}
\usepackage{fancyhdr}
\usepackage{needspace}
\usepackage{tikz}
\usetikzlibrary{arrows.meta,positioning,shapes.geometric,calc,fit,backgrounds}
\usepackage{hyperref}
\hypersetup{colorlinks=true,linkcolor=blu,citecolor=blu,urlcolor=blu}

\geometry{margin=2.5cm}

% --- colori coerenti con quelli delle figure (palette Okabe-Ito)
\definecolor{blu}{HTML}{0072B2}
\definecolor{arancio}{HTML}{E69F00}
\definecolor{verdes}{HTML}{009E73}
\definecolor{rossos}{HTML}{D55E00}
\definecolor{violas}{HTML}{CC79A7}
\definecolor{grigetto}{HTML}{E8E8E8}

\captionsetup{font=small,labelfont=bf,width=0.92\textwidth}

\pagestyle{fancy}
\fancyhf{}
\fancyhead[L]{\small\itshape\intestazione}
\fancyhead[R]{\small\thepage}
\renewcommand{\headrulewidth}{0.3pt}

% --- notazione
\newcommand{\ket}[1]{\lvert #1\rangle}
\newcommand{\bra}[1]{\langle #1 \rvert}
\newcommand{\media}[1]{\langle #1 \rangle}
\newcommand{\vs}[1]{\mathbf{s}_{#1}}

% --- riquadro "in una riga": il messaggio da portare a casa
\newcommand{\messaggio}[1]{%
  \begin{center}
  \begin{minipage}{0.93\textwidth}
  \setlength{\fboxsep}{9pt}%
  \colorbox{grigetto}{\begin{minipage}{\dimexpr\textwidth-2\fboxsep}
  \small\textbf{In una riga.}\enspace #1
  \end{minipage}}
  \end{minipage}
  \end{center}
}

% --- riquadro "come lo sappiamo": la verifica numerica dietro un'affermazione
\newcommand{\verifica}[1]{%
  \begin{center}
  \begin{minipage}{0.93\textwidth}
  \small\textcolor{verdes}{\rule{2.2em}{1.1pt}}\enspace
  \textcolor{verdes}{\textbf{Come lo sappiamo.}}\enspace #1
  \end{minipage}
  \end{center}
  \medskip
}
 subito dopo \documentclass.
% =====================================================================
\usepackage[T1]{fontenc}
\usepackage[utf8]{inputenc}
\usepackage[main=italian,provide=*]{babel}
\usepackage{amsmath,amssymb,mathtools}
\usepackage{graphicx}
\usepackage{booktabs}
\usepackage{colortbl}
\usepackage{xcolor}
\usepackage{geometry}
\usepackage{placeins}
\usepackage{caption}
\usepackage{enumitem}
\usepackage{fancyhdr}
\usepackage{needspace}
\usepackage{tikz}
\usetikzlibrary{arrows.meta,positioning,shapes.geometric,calc,fit,backgrounds}
\usepackage{hyperref}
\hypersetup{colorlinks=true,linkcolor=blu,citecolor=blu,urlcolor=blu}

\geometry{margin=2.5cm}

% --- colori coerenti con quelli delle figure (palette Okabe-Ito)
\definecolor{blu}{HTML}{0072B2}
\definecolor{arancio}{HTML}{E69F00}
\definecolor{verdes}{HTML}{009E73}
\definecolor{rossos}{HTML}{D55E00}
\definecolor{violas}{HTML}{CC79A7}
\definecolor{grigetto}{HTML}{E8E8E8}

\captionsetup{font=small,labelfont=bf,width=0.92\textwidth}

\pagestyle{fancy}
\fancyhf{}
\fancyhead[L]{\small\itshape\intestazione}
\fancyhead[R]{\small\thepage}
\renewcommand{\headrulewidth}{0.3pt}

% --- notazione
\newcommand{\ket}[1]{\lvert #1\rangle}
\newcommand{\bra}[1]{\langle #1 \rvert}
\newcommand{\media}[1]{\langle #1 \rangle}
\newcommand{\vs}[1]{\mathbf{s}_{#1}}

% --- riquadro "in una riga": il messaggio da portare a casa
\newcommand{\messaggio}[1]{%
  \begin{center}
  \begin{minipage}{0.93\textwidth}
  \setlength{\fboxsep}{9pt}%
  \colorbox{grigetto}{\begin{minipage}{\dimexpr\textwidth-2\fboxsep}
  \small\textbf{In una riga.}\enspace #1
  \end{minipage}}
  \end{minipage}
  \end{center}
}

% --- riquadro "come lo sappiamo": la verifica numerica dietro un'affermazione
\newcommand{\verifica}[1]{%
  \begin{center}
  \begin{minipage}{0.93\textwidth}
  \small\textcolor{verdes}{\rule{2.2em}{1.1pt}}\enspace
  \textcolor{verdes}{\textbf{Come lo sappiamo.}}\enspace #1
  \end{minipage}
  \end{center}
  \medskip
}

\providecommand{\Tr}{\operatorname{Tr}}

\title{\bfseries Punti di lavoro e ansatz nel dimero e nel trimero ad anello\\[2pt]
\large Riepilogo esaustivo --- dimero, trimero ad anello, estensioni facoltative}
\author{Samuele Schiavi \\ \small Tesi triennale in Fisica, Università di Parma
--- Relatore: Prof.\ Alessandro Chiesa}
\date{}

\begin{document}
\maketitle
\thispagestyle{fancy}

\begin{abstract}
\noindent
Questo documento raccoglie, per il dimero e per il trimero ad anello, tutti i
punti di lavoro $(b,D)$ usati nel progetto, gli ansatz variazionali impiegati
a ciascuno stadio, e lo stato delle due estensioni facoltative della pipeline
di rumore (ottimizzazione VQE noise-aware, readout asimmetrico). Include la
verifica appena condotta (``Passo 0'') sull'utilizzabilità del punto $R_0$
dell'anello come secondo punto di lavoro VQE, con la tabella completa dei
risultati. Nota di scope: il filone \emph{catena aperta} è stato messo da
parte per decisione esplicita e non compare qui se non per contestualizzare
la biforcazione fra ansatz $W$ e RBS.
\end{abstract}

\tableofcontents

% =====================================================================
\section{Il dimero: due punti di lavoro, non uno}
% =====================================================================

Il dimero usa lo stesso ansatz variazionale (\textbf{PMA-2q$\cdot$3}, 3
parametri, 2 CNOT) su \emph{due} punti di lavoro fisici distinti, entrambi
validati come stati fondamentali VQE genuini con $\mathcal F=1$ a precisione
di macchina.

\begin{table}[htbp]
\centering
\caption{Punti di lavoro del dimero.}
\label{tab:dimero-punti}
\begin{tabular}{lccl}
\toprule
Nome & $b/J$ & $D/J$ & Ruolo \\
\midrule
\textbf{test 2} & $0.35$ & $0.80$ & punto principale: VQE $\to$ ground state
$\to$ correlazioni \\
\textbf{$R_1$} & $-0.18$ & $1.00$ & nato per la dimostrazione Trotter
(dinamica non monocromatica);\\
& & & validato \emph{a posteriori} anche come punto VQE genuino \\
\bottomrule
\end{tabular}
\end{table}

\messaggio{$R_1$ non sostituisce ``test 2'': è stato aggiunto come secondo
punto di lavoro VQE, per verificare che i risultati delle estensioni non
fossero un caso isolato di un singolo punto.}

\subsection{Baseline della pipeline di rumore (cinque stadi), punto principale}

Al punto ``test 2'', calibrazione di riferimento \texttt{ibm\_torino}
(arXiv:2504.15187):
\begin{itemize}[nosep]
  \item $N^*=8$ sulla fedeltà di Trotter;
  \item $N^*=2$ sul modulo del correlatore dinamico $C_{11}^{yz}(t=2)$
  (valore corretto dopo revisione: $C_{21}^{xx}$, usato inizialmente, è
  risultato l'estremo più fragile su tutte le 36 combinazioni sotto
  rumore, sostituito con $C_{11}^{yz}$ --- vedi nota a fine sezione);
  \item $N^*$ dimostrato \textbf{analiticamente indipendente} da
  $p_\text{readout}$ nel caso simmetrico (la correzione di readout è un
  fattore moltiplicativo costante in $N$, non può spostare un massimo).
\end{itemize}

\begin{center}
\begin{minipage}{0.93\textwidth}
\small\textcolor{rossos}{\rule{2.2em}{1.1pt}}\enspace
\textcolor{rossos}{\textbf{Correzione (10 settembre 2026).}}\enspace
$C_{21}^{xx}$ è risultato, con uno scan sistematico sulle 36 combinazioni
sotto rumore, l'estremo più fragile in assoluto ($N^*$ più alto, segnale
più basso di tutte) --- non solo ``non ottimale'' come già segnalato nel
caso ideale. Sostituito con $C_{11}^{yz}$ (già usato come esempio ``ricco''
in Parte 1, verificato rappresentativo nella nuova distribuzione). Le
conclusioni qualitative delle estensioni (invarianza sotto riottimizzazione,
25/25 sulla griglia, stress test readout fino a $50\times$) restano confermate
con il nuovo correlatore --- cambiano solo i valori numerici di $N^*$. Nello
stesso giro, il readout è stato reso genuinamente campionato (shot finiti +
\texttt{ReadoutError} vero), non più formula analitica: stessa posizione di
$N^*$ confermata. \textbf{Nota aggiuntiva}: $|C(N^*)|$ sovrastima il
correlatore fisico vero di un fattore $\sim2\times$ al punto usato in tutta
la pipeline --- colpa dell'errore di Trotter a $N$ piccolo, non del rumore;
$N^*$ massimizza il segnale misurabile, non l'accuratezza della stima. Il
valore corretto di $N^*$ per $R_1$ (dopo la stessa sostituzione) non è
ancora noto a questo documento --- il valore $3\to3$ in Tabella 2 sotto è
quindi il numero \emph{precedente la correzione}, da verificare.
\end{minipage}
\end{center}
\medskip

% =====================================================================
\subsection{Estensione 1 --- VQE noise-aware, su entrambi i punti}
% =====================================================================

Due strategie a confronto: riuso dei parametri ideali $\theta^*_\text{ideale}$
su un simulatore a operatore densità, contro una riottimizzazione con il
rumore acceso fin dall'inizio della valutazione dell'energia
($\theta^*_\text{noise-aware}$), stesso schema (multistart COBYLA + polish
L-BFGS-B).

\textbf{Argomento teorico.} Per un canale depolarizzante isotropo,
$$E_\text{rumoroso}(\theta) = (1-\lambda)\,E_\text{ideale}(\theta) +
\lambda\,\frac{\Tr[H]}{d},$$
trasformazione \emph{affine} dell'energia ideale, indipendente da $\theta$ ---
vale per qualunque $\lambda$ e qualunque punto di lavoro. Questo non
garantisce da solo l'invarianza di $N^*$ su metriche diverse dall'energia
(fedeltà, modulo di un correlatore): va verificato direttamente.

\begin{table}[htbp]
\centering
\caption{Estensione VQE noise-aware sul dimero, confronto sui due punti di
lavoro. Riga $N^*$ correlatore: valori \emph{precedenti} la correzione del
correlatore (10 settembre) --- l'invarianza qualitativa (riuso $=$
noise-aware) resta confermata col nuovo correlatore $C_{11}^{yz}$, i valori
numerici esatti post-correzione non sono ancora disponibili a questo
documento (noto solo per test 2: $N^*=2$, non più $5$).}
\label{tab:dimero-noiseaware}
\begin{tabular}{lcc}
\toprule
Quantità & test 2 & $R_1$ \\
\midrule
$\Delta E = E_\text{noise-aware} - E_\text{riuso}$ & $-1.16\times10^{-7}$ & $-2.56\times10^{-8}$ \\
$\Delta \mathcal F$ & $+1.76\times10^{-8}$ & $+3.96\times10^{-9}$ \\
$N^*$ Trotter (riuso $\to$ noise-aware) & $8\to8$ & $8\to8$ \\
$N^*$ correlatore (riuso $\to$ noise-aware, \emph{pre-correzione}) & $5\to5$ & $3\to3$ \\
\bottomrule
\end{tabular}
\end{table}

\verifica{Argomento di continuità reso quantitativo:
$|\Delta f|\le\|\nabla f(\theta^*_\text{ideale})\|\,\|\Delta\theta\|+O(\|\Delta\theta\|^2)$.
La previsione lineare riproduce lo scostamento osservato entro lo
\textbf{0.1\%} su entrambe le metriche (fedeltà di Trotter: predetto
$4.150\times10^{-5}$ vs osservato $4.147\times10^{-5}$; correlatore: predetto
$-6.794\times10^{-5}$ vs osservato $-6.792\times10^{-5}$). Scan completo sulla
griglia del Passo 5 (7 valori di $\varepsilon_{2q}$, 6 di $\varepsilon_{1q}$
per la fedeltà; 5, 4, 3 valori di $\varepsilon_{2q},\varepsilon_{1q},
p_\text{readout}$ per il correlatore): \textbf{25/25 punti con $N^*$
invariato}.}

\textbf{Robustezza del multistart} ($R=6$, come nel Passo 2 originale): senza
seme, può restare intrappolato in un minimo locale peggiore
($E\approx-3.4753$ contro il vero $-3.50164605$). Risolto includendo
$\theta^*_\text{ideale}$ come uno dei punti di partenza; verificato
indipendentemente che non è un artefatto del seme, un multistart più ampio
($R=20$, nessun seme) converge allo stesso ottimo rumoroso entro
$2.6\times10^{-6}$.

\textbf{Esito}: nessun vantaggio misurabile dalla riottimizzazione, su
entrambi i punti, entrambe le metriche, tutta la griglia esplorata.

% =====================================================================
\subsection{Estensione 2 --- readout asimmetrico}
% =====================================================================

Formula generale di correzione per $p_{01}\neq p_{10}$:
$$\langle Z\rangle_\text{readout}=\alpha\langle Z\rangle_\text{ideale}+\beta,
\qquad \alpha = 1-p_{01}-p_{10},\ \ \beta = p_{10}-p_{01},$$
che per il correlatore complesso diventa una traslazione costante nel piano
complesso, $C_\text{readout}(N)=\alpha\,C_\text{ideale}(N)+\beta(1+i)$ ---
\textbf{non} garantita a priori preservare l'invarianza di
$N^*=\arg\max_N|C(N)|$ dimostrata per il caso simmetrico (somma e modulo non
commutano).

\textbf{Valori usati}: rapporto illustrativo $3\times$ ($p_{01}=0.0115$,
$p_{10}=0.0345$, stessa media $p_\text{readout}=2.3\times10^{-2}$ già in uso),
scelto confrontando split reali su chip IBM comparabili (IBM-Quito
$3.5\times$, esempio documentazione IBM $3.1\times$, ibmq\_montreal
$2.0\times$ --- nessuno specifico per \texttt{ibm\_torino} risultava
citabile).

\verifica{$N^*$ non si sposta (verificato $N^*=5$ con il correlatore
originale $C_{21}^{xx}$; valore atteso $N^*=2$ dopo la sostituzione con
$C_{11}^{yz}$, non ancora riverificato su questo test specifico), né al
punto illustrativo né in uno stress test fino a rapporti
$p_{10}/p_{01}=50\times$ (ben oltre il range realistico osservato su
hardware). Cross-check Monte Carlo con \texttt{ReadoutError} vero
(asimmetrico) conferma la formula analitica entro l'errore statistico
atteso.}

% =====================================================================
\section{Il trimero ad anello: due punti di lavoro con ruoli diversi}
% =====================================================================

\subsection{Biforcazione dell'ansatz: RBS e $W$}

A $D=0$ le due famiglie di ansatz sono state sviluppate \textbf{in
parallelo}, per entrambe le topologie (anello e catena): un ansatz PMA con
blocco RBS (produzione principale, più efficiente in questo regime) e uno
speculare con blocco $W$ (mirror dell'ansatz ``originale'' del dimero).

Sotto il termine DM (Fase 5), il confronto sistematico condotto nel notebook
dedicato dell'anello (15 ansatz: 10 RBS, 5 $W$, 60 restart, ottimizzazione
diretta della fidelity) ha dato un esito netto:
\begin{align*}
\text{RBS-2q (1 giro)} &\ \to\ \mathcal F\approx0.9999\quad\text{(tetto strutturale)},\\
W\text{-2q.6} &\ \to\ \mathcal F=1\quad\text{esatto}.
\end{align*}

Da quel momento \textbf{\texttt{W-2q.6}} (6 parametri, blocco
CNOT--$R_y(\theta)$--CNOT, un solo giro sui tre legami $(1,2),(2,3),(3,1)$) è
diventato l'ansatz canonico unico per tutto ciò che segue nella pipeline
dell'anello --- circuito dei correlatori, riferimento Trotter, VQE+DM sotto
rumore. RBS non è stato portato avanti oltre il notebook di confronto.

\begin{center}
\begin{minipage}{0.93\textwidth}
\small\textcolor{rossos}{\rule{2.2em}{1.1pt}}\enspace
\textcolor{rossos}{\textbf{Precisazione trovata a posteriori.}}\enspace
La frase ``RBS-2q ha un tetto strutturale vero'' vale solo per la variante a
\emph{un giro} effettivamente testata: nello sweep sistematico, \texttt{RBS-2qC.K2}
(due giri) raggiunge anch'esso $\mathcal F=1$ esatto sotto DM. Nessuna azione
correttiva sui file consegnati --- la conclusione operativa (\texttt{W-2q.6}
canonico) resta valida, solo la portata della frase è più stretta di una
lettura letterale.
\end{minipage}
\end{center}
\medskip

\subsection{I due punti di lavoro dell'anello}

A differenza del dimero, l'anello usa due punti \textbf{nati per scopi
diversi} in Parte 1, non entrambi originariamente pensati come punti VQE.

\begin{table}[htbp]
\centering
\caption{Punti di lavoro del trimero ad anello ($J=1$, $J'=0.4$ in entrambi).}
\label{tab:anello-punti}
\begin{tabular}{lccl}
\toprule
Nome & $b$ & $D$ & Ruolo originario \\
\midrule
\textbf{Scenario A} & $b_c=2.4$ & $0.15$ & punto VQE (Fase 5), gap $E_1-E_0=0.2547$ \\
\textbf{$R_0$} & $0.05$ & $1.93$ & dimostrazione Trotter (dinamica non
monocromatica),\\
& & & \emph{dichiarato provvisorio}, gap $E_1-E_0=0.0393$ \\
\bottomrule
\end{tabular}
\end{table}

\messaggio{Il gap di $R_0$ è $6.5\times$ più piccolo di quello dello Scenario
A --- $R_0$ è vicino alla quasi-degenerazione di Kramers che il modello ha
esattamente a $b=0$ (per qualunque $D$, con un numero dispari di spin
$1/2$).}

\subsection{Baseline della pipeline di rumore (cinque stadi), entrambi gli scenari}

A differenza del dimero, la pipeline di rumore dell'anello tratta
\textbf{entrambi} gli scenari fin dallo Stadio 3, non uno scelto a scapito
dell'altro.

\begin{table}[htbp]
\centering
\caption{Baseline dei cinque stadi della pipeline di rumore, trimero ad anello.}
\label{tab:anello-stadi}
\begin{tabular}{lll}
\toprule
Stadio & Contenuto & Esito \\
\midrule
1 & modello di rumore & calibrazione \texttt{ibm\_torino} riusata dal dimero \\
2 & VQE+DM sotto rumore & $E=-5.37687489$ (ideale $-5.53437002$), $1-\mathcal F=2.98\%$\\
  & & (doppio del dimero, $1.50\%$: 6 CNOT contro 2) \\
3 & Trotter sotto rumore & $N^*_A=3$, $N^*_B=9$ (dopo correzione griglia) \\
4 & correlatori sotto rumore & $N^*=3$ per $C_{21}^{xx}$ (coincide con $N^*_A$) \\
5 & scan parametri di rumore & $N^*$ monot.\ non-cresc.\ in $\varepsilon_{1q},\varepsilon_{2q}$;\\
  & & invarianza da $p_\text{readout}\in\{0,0.05,0.20\}$ confermata \\
\bottomrule
\end{tabular}
\end{table}

\begin{center}
\begin{minipage}{0.93\textwidth}
\small\textcolor{arancio}{\rule{2.2em}{1.1pt}}\enspace
\textcolor{arancio}{\textbf{Correzione numerica non banale.}}\enspace
Una griglia $N$ troppo grezza nel primo giro dello Stadio 3 aveva dato
$N^*_A=2$, $N^*_B=10$ invece dei valori corretti $3$ e $9$ --- scoperta
costruendo lo Stadio 5, propagata retroattivamente. Conseguenza: $N^*_A=3$
\textbf{coincide} esattamente con l'$N^*=3$ del correlatore (Stadio 4), prima
sembravano tre valori distinti.
\end{minipage}
\end{center}
\medskip

\textbf{Verifica di minimizzazione dei gate} (richiesta esplicita del
relatore, ``farei prima la transpilazione, aiuterà a ridurre i gate''):
fattore $2\times$ sul passo esterno di Trotter ($30\to15$ CNOT, livelli 0/1
contro 2/3); \textbf{nessun beneficio} sull'ansatz VQE a nessun livello (6
CNOT costanti, struttura senza ridondanza da fondere).

% =====================================================================
\section{Passo 0 --- verifica che $R_0$ sia un punto VQE valido}
% =====================================================================

Prerequisito individuato prima di poter riusare $R_0$ come secondo punto di
lavoro nell'estensione VQE noise-aware (per analogia col ruolo di $R_1$ sul
dimero): $R_0$ non era mai stato validato come stato fondamentale VQE, solo
come punto di dimostrazione della dinamica Trotter da $|000\rangle$.

\textbf{Preoccupazione motivata}: $b=0.05$ è vicino a $b=0$, dove il modello
ha una degenerazione di Kramers esatta per qualunque $D$ (numero dispari di
spin $1/2$) --- un paesaggio di ottimizzazione vicino a una quasi-degenerazione
può presentare direzioni piatte o convergenza instabile a seconda del seed.

\subsection{Metodo}

VQE con ansatz \texttt{W-2q.6} al punto $R_0$ ($J=1,J'=0.4,b=0.05,D=1.93$,
Opzione DM B), stesso schema multistart COBYLA (\texttt{maxiter}$=2000$) +
polish L-BFGS-B già usato per lo Scenario A, ripetuto su 5 seed indipendenti
con $R=12$ riavvii ciascuno, più un controllo di robustezza a $R=24$ su un
seed indipendente.

\subsection{Risultati}

\begin{table}[htbp]
\centering
\caption{Passo 0: VQE su \texttt{W-2q.6} al punto $R_0$, multistart su 5 seed
indipendenti più un controllo di robustezza. $E_0=-6.2971308468$ (esatto),
gap $=0.039274$.}
\label{tab:passo0}
\begin{tabular}{ccccc}
\toprule
seed & $R$ & $E_\text{vqe}$ & $\Delta E$ & $\mathcal F$ \\
\midrule
0   & 12 & $-6.2971308468$ & $9.299\times10^{-13}$ & $0.999999999998$ \\
1   & 12 & $-6.2971308468$ & $5.063\times10^{-14}$ & $0.999999999999$ \\
2   & 12 & $-6.2971308468$ & $1.368\times10^{-13}$ & $0.999999999997$ \\
3   & 12 & $-6.2971308468$ & $1.757\times10^{-12}$ & $0.999999999956$ \\
4   & 12 & $-6.2971308468$ & $8.260\times10^{-14}$ & $0.999999999998$ \\
100 & 24 & $-6.2971308468$ & $8.882\times10^{-14}$ & $0.999999999998$ \\
\bottomrule
\end{tabular}
\end{table}

Tutti e 6 i run convergono allo \textbf{stesso ottimo globale}: nessuna
traccia di minimi locali spuri. I parametri ottimi differiscono
numericamente da seed a seed (dispersione alta componente per componente),
ma a parità di energia e fidelity --- coerente con simmetrie interne
dell'ansatz (parametrizzazioni multiple equivalenti dello stesso stato
fisico), non con ottimi distinti.

\verifica{5 seed a $R=12$ più un controllo indipendente a $R=24$, tutti
convergenti allo stesso ottimo entro $\Delta E\lesssim2\times10^{-12}$ e
$\mathcal F\geq0.999999999956$ --- lo stesso standard di verifica già
applicato allo Scenario A (12/12 multistart convergenti).}

\messaggio{Il gap più piccolo di $R_0$ non si traduce in un problema di
ottimizzazione: il multistart converge in modo pulito su tutti i tentativi.
$R_0$ è un secondo punto di lavoro VQE genuino, utilizzabile direttamente
nell'estensione VQE noise-aware --- nessuno scan alternativo necessario.}

% =====================================================================
\section{Stato delle estensioni sull'anello e piano di lavoro}
% =====================================================================

\begin{table}[htbp]
\centering
\caption{Stato delle due estensioni facoltative, trimero ad anello.}
\label{tab:anello-estensioni-stato}
\begin{tabular}{lll}
\toprule
Estensione & Fattibilità & Stato \\
\midrule
1. VQE noise-aware & con adattamenti & Passo 0 chiuso; Passo 1 \\
   & & (solo Scenario A) \textbf{chiuso} \\
2. Readout asimmetrico & senza modifiche concettuali & \textbf{da fare} \\
   & & (solo l'ancilla è misurata nel circuito \\
   & & a 4 qubit --- stesso schema del dimero) \\
\bottomrule
\end{tabular}
\end{table}

\begin{center}
\begin{minipage}{0.93\textwidth}
\small\textcolor{rossos}{\rule{2.2em}{1.1pt}}\enspace
\textcolor{rossos}{\textbf{Scoping del Passo 1, deciso dopo il Passo 0.}}\enspace
Verificato nel codice (\texttt{trotter\_rumoroso\_trimero\_anello.py},
\texttt{correlatori\_rumorosi\_trimero\_anello.py}) che $R_0$, nella
pipeline attuale, non ha alcuna preparazione VQE da riottimizzare: il
Trotter a $R_0$ parte da $\ket{000}$ (non dal ground state VQE), e il
modulo dei correlatori dichiara esplicitamente di non usare mai $R_0$.
``VQE noise-aware'' presuppone una preparazione VQE da riottimizzare sotto
rumore --- condizione soddisfatta solo dallo Scenario A. Si procede quindi
con \textbf{Scenario A soltanto}, su entrambe le metriche (fedeltà di
Trotter e modulo del correlatore, baseline $N^*_A=3$ ed $N^*=3$) ---
analogo diretto di ``test 2'' sul dimero. $R_0$ resta quello già definito
in Parte 1 (dimostrazione Trotter da $\ket{000}$), non toccato da questa
estensione; l'uso del punto VQE $R_0$ scoperto nel Passo 0 resta
disponibile per un'eventuale estensione futura, non inclusa qui.
\end{minipage}
\end{center}
\medskip

\textbf{Ostacoli concreti identificati per l'estensione 1} (rispetto al
dimero): paesaggio a 6 parametri invece di 3; costo per valutazione
$\sim4$--$10\times$ maggiore (operatore densità $8\times8$ contro $4\times4$,
6 CNOT contro 2); disciplina di transpilazione da replicare da zero.

\subsection{Passo 1 (VQE noise-aware, Scenario A) --- chiuso: un comportamento nuovo, non visto sul dimero}

Primo multistart ($R=12$ + seme $\theta^*_\text{ideale}$, rumore di
riferimento \texttt{ibm\_torino}): $\Delta E\approx-0.046$ (tre ordini di
grandezza più grande del $\Delta E\sim10^{-7}$ del dimero), ma
$\Delta\mathcal F<0$ --- la fedeltà \textbf{peggiora}, non migliora.

\begin{center}
\begin{minipage}{0.93\textwidth}
\small\textcolor{arancio}{\rule{2.2em}{1.1pt}}\enspace
\textcolor{arancio}{\textbf{Meccanismo isolato: sfruttamento del transpilatore.}}\enspace
I parametri trovati hanno due componenti quasi esatte a $0$ e $\pi$
($\theta_1/\pi\approx0.99999637$, $\theta_2/\pi\approx0.00889$). A questi
valori speciali il blocco CNOT--$R_y(\theta)$--CNOT si semplifica, e il
transpilatore (livello 3, quello di produzione) elimina 2 CNOT su 6.
Verificato su più seed ($R=12$, seed 0--4, più un controllo a $R=24$): il
comportamento a livello 3 è \textbf{fortemente instabile}
($E\in[-5.423,-5.377]$, $\mathcal F\in[0.70,0.96]$, CNOT oscilla fra 4 e 6
a seconda del seed) --- l'ottimizzatore sta in parte sfruttando un effetto
discontinuo del transpilatore, non solo esplorando la fisica del problema.
\end{minipage}
\end{center}
\medskip

\textbf{Isolato il meccanismo con un controllo a struttura fissa}
(\texttt{optimization\_level=0}, verificato dare \emph{sempre} lo stesso
conteggio di gate --- 18 rz + 12 sx + 6 cx --- qualunque $\theta$, anche ai
valori speciali sopra):

\begin{table}[htbp]
\centering
\caption{VQE noise-aware, Scenario A: struttura di circuito variabile
(produzione) contro struttura fissa. 3/3 seed a struttura fissa
convergono identici; a struttura variabile il risultato dipende fortemente
dal seed.}
\label{tab:anello-noiseaware-preliminare}
\begin{tabular}{lccc}
\toprule
& $E$ & $\mathcal F$ & CNOT \\
\midrule
Riuso ideale (livello 3) & $-5.37687$ & $0.97021$ & 6 \\
Noise-aware, livello 3, seed 0 & $-5.42250$ & $0.96383$ & 4 \\
Noise-aware, livello 3, seed 1 & $-5.42279$ & $0.90994$ & 4 \\
Noise-aware, livello 3, seed 2 & $-5.37707$ & $0.69969$ & 4 \\
Noise-aware, livello 3, seed 3 & $-5.37779$ & $0.93575$ & 6 \\
Noise-aware, livello 3, seed 4 & $-5.37809$ & $0.96246$ & 6 \\
\midrule
Riuso ideale (livello 0, struttura fissa) & $-5.36232$ & $0.96728$ & 6 \\
Noise-aware, livello 0, seed 0/1/2 & $-5.36649$ & $0.95154$ & 6 \\
\bottomrule
\end{tabular}
\end{table}

A struttura fissa il paesaggio è ben comportato, e mostra un effetto fisico
piccolo ma reale, di segno inatteso: $\Delta E=-0.00416$ (ancora tre ordini
di grandezza più grande del dimero) accompagnato da $\Delta\mathcal
F=-0.0158$ --- energia migliore, fedeltà peggiore, un compromesso che il
dimero non mostrava affatto.

\messaggio{Sull'anello, a differenza del dimero, riottimizzare sotto
rumore non è innocuo: porta a un'energia più bassa ma a una fedeltà più
bassa. Il comportamento erratico a struttura variabile (livello di
produzione) è un fenomeno distinto e reale --- il transpilatore introduce
un incentivo implicito a raggiungere configurazioni particolari di
$\theta$ --- da documentare a parte, non da scartare come errore.}

\subsubsection*{Controlli $N^*$ a valle --- completati}

\begin{table}[htbp]
\centering
\caption{$N^*$ a valle, Scenario A, riuso ideale contro noise-aware, sulle
due metodologie. Baseline noto: $N^*_A=3$ (Trotter), $N^*=3$ (correlatore).}
\label{tab:anello-Nstar-passo1}
\begin{tabular}{lcc}
\toprule
Metodologia & $N^*$ Trotter & $N^*$ correlatore \\
\midrule
Struttura fissa (livello 0, risultato fisico) & $3\to3$ & $3\to3$ \\
Struttura di produzione (livello 3, punto anomalo) & $3\to\mathbf{4}$ & $3\to3$ \\
\bottomrule
\end{tabular}
\end{table}

A struttura fissa, $N^*$ resta invariato su \textbf{entrambe} le metriche
--- stesso schema qualitativo del dimero, nonostante il segno diverso
dell'effetto a monte (energia migliore, fedeltà peggiore anziché
entrambe quasi invariate). Sul punto anomalo (struttura variabile,
produzione reale), $N^*$ \textbf{si sposta} per la fedeltà di Trotter
($F(3)=0.700335$ contro $F(4)=0.693018$, scarto piccolo ma il massimo
cambia posizione) mentre resta invariato per il correlatore ---
un'asimmetria fra le due metriche mai osservata sul dimero, dove
l'invarianza reggeva sempre su entrambe.

\messaggio{Passo 1 chiuso. Risultato centrale: isolando l'effetto fisico a
struttura di circuito fissa, la riottimizzazione sotto rumore non sposta
$N^*$ sull'anello, esattamente come sul dimero --- ma la pipeline di
produzione (struttura variabile) può occasionalmente farlo su una sola
delle due metriche, un comportamento specifico dell'ansatz a più
parametri, da riportare come nota metodologica a sé.}

\subsection{Verifica di non-regressione: fragilità del correlatore sull'anello}

A seguito della correzione ricevuta dalla revisione del dimero
($C_{21}^{xx}$ risultato l'estremo più fragile su 36 combinazioni sotto
rumore, sostituito da $C_{11}^{yz}$, $N^*=5\to2$): verificato se lo stesso
rischio riguarda il baseline dell'anello ($C_{21}^{xx}$, $N^*=3$, Stadio 4),
mai controllato prima con questo criterio.

\begin{table}[htbp]
\centering
\caption{Scan sotto rumore sulle 81 combinazioni, trimero ad anello
(\texttt{ibm\_torino}, griglia $N\in\{1,2,3,4,5,7,10,14,20\}$).}
\label{tab:anello-scan81}
\begin{tabular}{lccc}
\toprule
& $N^*$ & $|C(N^*)|$ & posizione per segnale (su 81) \\
\midrule
$C_{21}^{xx}$ (baseline Stadio 4) & 3 & 0.5048 & 64° (sopra la mediana) \\
$C_{33}^{xz}$ (il più fragile) & 2 & 0.0461 & 1° \\
$C_{33}^{yz}$ (il più robusto) & 1 & 0.7255 & 81° \\
\bottomrule
\end{tabular}
\end{table}

Distribuzione di $N^*$ sulle 81 combinazioni: $N^*{=}1$: 40, $N^*{=}2$: 18,
$N^*{=}3$: 21, $N^*{=}4$: 2 (solo $C_{23}^{yz}$ e $C_{13}^{zz}$). Il valore
$N^*=3$ di $C_{21}^{xx}$ coincide esattamente col baseline già noto ed è
condiviso da altre 20 combinazioni --- non un caso isolato.

\messaggio{Nessuna correzione necessaria sul baseline dell'anello: a
differenza del dimero, $C_{21}^{xx}$ è una scelta rappresentativa anche
sotto rumore, non solo nel caso ideale.}

\textbf{Nessun ostacolo concettuale per l'estensione 2}: verificato
direttamente dal codice (\texttt{circuito\_correlazioni\_trimero\_anello.py})
che il circuito misura un solo bit classico, sempre e solo l'ancilla
(\texttt{qc.measure(ANCILLA, 0)}) --- la correzione di readout asimmetrico si
applica alla stessa, singola matrice di confusione indipendentemente dal
numero di qubit di registro (3 sull'anello contro 2 sul dimero). Da
allineare, quando affrontata, alla metodologia ora in uso sul dimero
(readout genuinamente campionato, non solo formula analitica).

% =====================================================================
\section{Riepilogo comparativo finale}
% =====================================================================

\begin{table}[htbp]
\centering
\caption{Confronto dimero / trimero ad anello: ansatz, dimensione, costo.}
\label{tab:confronto-finale}
\begin{tabular}{lcc}
\toprule
& Dimero & Trimero ad anello \\
\midrule
Ansatz canonico sotto DM & PMA-2q$\cdot$3 & $W$-2q.6 \\
Parametri & 3 & 6 \\
CNOT (ansatz) & 2 & 6 \\
Qubit di registro & 2 & 3 \\
Qubit circuito correlatori (+ancilla) & 3 & 4 \\
Punti di lavoro usati & test 2, $R_1$ & Scenario A, $R_0$ \\
$N^*$ Trotter (rumore rif.) & 8 & $3$ (A) / $9$ ($R_0$) \\
$N^*$ correlatore (rumore rif.) & 2 (era 5, corretto) & 3 (verificato non fragile) \\
Estensione VQE noise-aware & chiusa & Passo 0 e Passo 1 chiusi \\
Estensione readout asimmetrico & chiusa$^\dagger$ & da fare \\
\bottomrule
\end{tabular}

\vspace{2pt}
{\footnotesize $^\dagger$ da riallineare alla metodologia ora in uso sul
dimero (readout genuinamente campionato, non solo formula analitica).}
\end{table}

\end{document}
